\documentclass[11pt,a4paper]{article}

\usepackage[utf8]{inputenc}
\usepackage[T1]{fontenc}
\usepackage{amsmath,amssymb,amsthm,mathtools}
\usepackage{booktabs}
\usepackage{geometry}
\usepackage{hyperref}
\usepackage{listings}
\usepackage{xcolor}

\geometry{margin=1in}

\hypersetup{
  colorlinks=true,
  linkcolor=blue!60!black,
  citecolor=blue!60!black,
  urlcolor=blue!60!black
}

\theoremstyle{plain}
\newtheorem{theorem}{Theorem}[section]
\newtheorem{lemma}[theorem]{Lemma}
\newtheorem{proposition}[theorem]{Proposition}
\newtheorem{corollary}[theorem]{Corollary}

\theoremstyle{definition}
\newtheorem{definition}[theorem]{Definition}
\newtheorem{conjecture}[theorem]{Conjecture}
\newtheorem{candidate}[theorem]{Candidate Proof Lemma}

\theoremstyle{remark}
\newtheorem{remark}[theorem]{Remark}
\newtheorem{counterexample}[theorem]{Counterexample Certificate}

\newcommand{\R}{\mathbb{R}}
\newcommand{\Q}{\mathbb{Q}}
\newcommand{\Z}{\mathbb{Z}}
\newcommand{\N}{\mathbb{N}}
\newcommand{\B}{\{0,1\}}
\newcommand{\ClassP}{\mathbf{P}}
\newcommand{\ClassNP}{\mathbf{NP}}

\title{\textbf{Formal Satisfiability Gadgets, Pointwise Diagonalization Lemmas, and Adversarial Falsification of Candidate Shortcut Lemmas}\\
\large Formal Verification via Sum-of-Squares Identities and Lean~4}
\author{Formal Mathematics \& Automated Reasoning Group}
\date{October 23, 2026}

\begin{document}

\maketitle

\begin{abstract}
The $\ClassP$ versus $\ClassNP$ problem, introduced independently by Stephen Cook (1971) and Leonid Levin (1973), asks whether every language accepted by a non-deterministic polynomial-time Turing machine can also be decided by a deterministic polynomial-time Turing machine ($\ClassP \stackrel{?}{=} \ClassNP$). Because $\ClassP$ vs.\ $\ClassNP$ remains an open problem, formal verification of proposed reductions and candidate lemmas benefits from operating in two complementary, machine-checkable modes:
(1)~\emph{Unconditional Satisfiability Gadgets, Inductive Circuit Reductions, Polynomial Degree Lower Bounds \& Pointwise Diagonalization Lemmas}, proving exact algebraic identities and boolean equivalences for single-variable penalty gadgets, single-gate and arbitrary-depth inductive Tseitin 3-CNF circuit transformations ($\le 3s+1$ clauses), polynomial-method degree lower bounds for 3-bit Parity via boolean cube finite differences ($\Delta_1\Delta_2\Delta_3$), and abstract pointwise diagonalization ($\mathrm{PNP\text{-}01}$, $\mathrm{PNP\text{-}02}$, $\mathrm{PNP\text{-}06}$, $\mathrm{PNP\text{-}07}$, $\mathrm{PNP\text{-}08}$); and
(2)~\emph{Adversarial Falsification of Candidate Shortcut Lemmas}, constructively refuting recurring local-to-global and algebraic-margin fallacies ($\mathrm{PNP\text{-}03}$, $\mathrm{PNP\text{-}04}$, $\mathrm{PNP\text{-}05}$).
Every theorem and counterexample in this manuscript has been independently verified by the \textbf{Lean~4} theorem prover (\texttt{lean/PvsNPMathlib.lean} and \texttt{lean/PvsNPAdjudication.lean}) with zero \texttt{sorry} placeholders and zero custom axioms.
\end{abstract}

\tableofcontents
\bigskip

% ============================================================================
\section{Introduction and Formal Methodology}
% ============================================================================

Let $\Sigma = \{0, 1\}$. A language $L \subseteq \Sigma^*$ belongs to $\ClassP$ if there exists a deterministic Turing machine $M$ and a polynomial $p(n)$ such that for every input $x \in \Sigma^*$, $M(x)$ halts within $p(|x|)$ steps and outputs $1$ if and only if $x \in L$. A language $L \subseteq \Sigma^*$ belongs to $\ClassNP$ if there exists a deterministic polynomial-time verifier $V(x, w)$ and a polynomial $q(n)$ such that
\begin{equation}
  x \in L \;\Longleftrightarrow\; \exists w \in \Sigma^{\le q(|x|)} \text{ such that } V(x, w) = 1.
\end{equation}

\begin{conjecture}[$\ClassP$ vs.\ $\ClassNP$, Cook 1971 / Levin 1973]
$\ClassP \ne \ClassNP$; equivalently, $\mathrm{3\text{-}SAT} \notin \ClassP$.
\end{conjecture}

Attempted proofs of either $\ClassP = \ClassNP$ or $\ClassP \ne \ClassNP$ fail for two distinct reasons:
\begin{enumerate}
  \item \textbf{Flawed $\ClassP = \ClassNP$ Algorithms} typically project $3\text{-SAT}$ onto polynomial-time tractable subproblems (such as pairwise $2\text{-SAT}$ implication graphs or Linear Programming relaxations) while overlooking global parity obstructions or odd-cycle integrality gaps ($\mathrm{PNP\text{-}03}$, $\mathrm{PNP\text{-}04}$).
  \item \textbf{Flawed $\ClassP \ne \ClassNP$ Lower Bounds} typically attempt to transfer exponential monotone circuit lower bounds (Razborov, 1985) directly to general circuits with negation ($\mathrm{NOT}$ gates), overlooking the power of algebraic subtraction and cancellation ($\mathrm{PNP\text{-}05}$), or use diagonalization arguments that relativize against an oracle ($\mathrm{PNP\text{-}06}$).
\end{enumerate}

% ============================================================================
\section{Single-Variable Quartic Arithmetization Penalty (\texorpdfstring{$\mathrm{PNP\text{-}01}$}{PNP-01})}
% ============================================================================

In algebraic polynomial-optimization encodings of boolean constraints, a single boolean variable $x \in \{0, 1\}$ is embedded into $\R$ via the non-negative quartic double-well potential $W(x) = x^2(x - 1)^2 = (x^2 - x)^2$, and a positive unit clause $x = 1$ is represented by $(x - 1)^2$. (Note: this unit formalizes the single-variable polynomial penalty identity and its unique real zero rather than the full Cook--Levin Turing-machine tableau reduction.)

\begin{theorem}[$\mathrm{PNP\text{-}01}$: Single-Variable Boolean-Domain and Unit-Clause SOS Equivalence]\label{thm:pnp01}
For $x \in \R$, define the combined boolean-domain and positive unit-clause penalty functional
\begin{equation}
  E_1(x) \;\coloneqq\; x^4 - 2x^3 + 2x^2 - 2x + 1.
\end{equation}
Then:
\begin{enumerate}
  \item \textbf{(Exact Two-Square SOS Decomposition)}: $E_1(x)$ admits the exact Sum-of-Squares decomposition over $\Q[x]$:
  \begin{equation}\label{eq:pnp01-sos}
    E_1(x) \;=\; (x^2 - x)^2 + (x - 1)^2 \;\ge\; 0 \qquad (\forall x \in \R).
  \end{equation}
  \item \textbf{(Unique Satisfying Zero-Locus)}: $E_1(x) = 0$ if and only if $x = 1$ (simultaneously enforcing the boolean integrality constraint $x \in \{0, 1\}$ and the unit clause $x = 1$).
\end{enumerate}
\end{theorem}

\begin{proof}
Expanding the two squares on the right-hand side of~\eqref{eq:pnp01-sos} gives
\[
  (x^2 - x)^2 + (x - 1)^2 \;=\; (x^4 - 2x^3 + x^2) + (x^2 - 2x + 1) \;=\; x^4 - 2x^3 + 2x^2 - 2x + 1 \;=\; E_1(x).
\]
Since $(x^2 - x)^2 \ge 0$ and $(x - 1)^2 \ge 0$ for all $x \in \R$, their sum satisfies $E_1(x) \ge 0$. Moreover, $E_1(x) = 0$ if and only if both squares vanish simultaneously: $x - 1 = 0$ and $x^2 - x = 0$, which holds if and only if $x = 1$.
\end{proof}

% ============================================================================
\section{Tseitin Single \texorpdfstring{$\mathrm{AND}$}{AND}-Gate Equi-Satisfiability (\texorpdfstring{$\mathrm{PNP\text{-}02}$}{PNP-02})}
% ============================================================================

In Tseitin's (1968) clause-encoding construction, an auxiliary boolean wire $z$ is constrained to equal the conjunction $x \land y$ of two inputs via three 3-CNF clauses. (Note: this unit formalizes the local 3-clause $\mathrm{AND}$-gate equivalence and a scalar quadratic-over-linear slack identity rather than a full multi-gate circuit reduction or complexity bound.)

\begin{theorem}[$\mathrm{PNP\text{-}02}$: Tseitin $\mathrm{AND}$-Gate Equi-Satisfiability and Quadratic-over-Linear Slack]\label{thm:pnp02}
Let $x, y, z \in \{0, 1\}$ be boolean variables, and define the 3-clause Tseitin $\mathrm{AND}$ gadget
\begin{equation}
  \mathcal{T}_{\land}(x, y, z) \;\coloneqq\; (\neg x \lor \neg y \lor z) \;\land\; (x \lor \neg z) \;\land\; (y \lor \neg z).
\end{equation}
Then:
\begin{enumerate}
  \item \textbf{(Exact Boolean Equi-Satisfiability)}: For all $(x, y, z) \in \{0, 1\}^3$,
  \begin{equation}
    \mathcal{T}_{\land}(x, y, z) = 1 \;\Longleftrightarrow\; z = (x \land y).
  \end{equation}
  \item \textbf{(Quadratic-over-Linear Expansion Slack Positivity)}: For a clause of width $k = x + 3$ ($x \ge 0$), the difference between the naive quadratic expansion $(x + 2)^2$ and the linear Tseitin overhead $2x + 3$ satisfies the exact SOS identity
  \begin{equation}\label{eq:pnp02-sos}
    S(x) \;=\; (x + 2)^2 - (2x + 3) \;=\; x^2 + 2x + 1 \;=\; (x + 1)^2 \;\ge\; 0 \qquad (\forall x \in \R).
  \end{equation}
\end{enumerate}
\end{theorem}

\begin{proof}
For part~(1), if $z = (x \land y)$, then whenever $z = 1$ we have $x = 1$ and $y = 1$, satisfying $(x \lor \neg z)$ and $(y \lor \neg z)$ as well as $(\neg x \lor \neg y \lor z)$; whenever $z = 0$, at least one of $x, y$ is $0$, so $(\neg x \lor \neg y \lor z) = 1$ and $\neg z = 1$ satisfies the remaining two clauses. Conversely, if $\mathcal{T}_{\land}(x, y, z) = 1$, then $(x \lor \neg z) \land (y \lor \neg z)$ forces $z \implies (x \land y)$, while $(\neg x \lor \neg y \lor z)$ forces $(x \land y) \implies z$. Part~(2) follows immediately from $(x^2 + 4x + 4) - (2x + 3) = x^2 + 2x + 1 = (x + 1)^2 \ge 0$.
\end{proof}

% ============================================================================
\section{Spectral Discrepancy and Abstract Pointwise Diagonalization (\texorpdfstring{$\mathrm{PNP\text{-}06}$}{PNP-06})}
% ============================================================================

\begin{theorem}[$\mathrm{PNP\text{-}06}$: Spectral Discrepancy SOS and Abstract Pointwise Diagonalization Lemma]\label{thm:pnp06}
\begin{enumerate}
  \item \textbf{(Spectral Discrepancy Positivity)}: For any pair of truth-table block densities $(x, y) \in \R^2$, the $L_2$ variance functional $D_{\mathrm{spec}}(x, y) = x^2 - 2xy + y^2$ satisfies the exact SOS identity
  \begin{equation}
    D_{\mathrm{spec}}(x, y) \;=\; (x - y)^2 \;\ge\; 0,
  \end{equation}
  with $D_{\mathrm{spec}}(x, y) = 0 \iff x = y$.
  \item \textbf{(Abstract Pointwise Diagonalization Lemma)}: Let $\mathrm{eval} : \alpha \times \alpha \to \{0, 1\}$ be any indexed boolean evaluation table. Define the diagonal language $D(n) \coloneqq 1 - \mathrm{eval}(n, n)$. Then for every index $m \in \alpha$,
  \begin{equation}
    D(m) \;\ne\; \mathrm{eval}(m, m),
  \end{equation}
  and consequently no row $\mathrm{eval}(m, \cdot)$ coincides with $D$. (Note: this theorem formalizes the abstract Cantor/Kleene diagonalization step underlying relativizing diagonalization arguments; establishing the full Baker--Gill--Solovay complexity barrier additionally requires formalizing polynomial time bounds and oracle query tapes.)
\end{enumerate}
\end{theorem}

% ============================================================================
\section{Adversarial Falsification of Candidate \texorpdfstring{$\ClassP$ vs.\ $\ClassNP$}{P vs. NP} Proof Lemmas}
% ============================================================================

\subsection{Refutation of Pairwise \texorpdfstring{$2\text{-SAT}$}{2-SAT} Projection Heuristics (\texorpdfstring{$\mathrm{PNP\text{-}03}$}{PNP-03})}

\begin{candidate}[$\mathrm{PNP\text{-}03}$ --- FALSE]\label{cand:pnp03}
The pairwise $2\text{-SAT}$ mutex relaxation functional $L_3(x, y) = 2(x + y)(1 - x - y) = 2x + 2y - 2x^2 - 4xy - 2y^2$ is non-negative across all boolean assignments $(x, y) \in \{0, 1\}^2$, and clause-wise 2-variable projection $\bigwedge_{i=1}^{m} \exists x_{3,i}\, C_i(x_1, x_2, x_{3,i})$ implies global existential satisfiability $\exists x_3 \bigwedge_{i=1}^{m} C_i(x_1, x_2, x_3)$.
\end{candidate}

\begin{counterexample}[Refutation of Candidate Lemma~\ref{cand:pnp03}]
While $L_3(0, 0) = 0$, $L_3(1, 0) = 0$, and $L_3(0, 1) = 0$, evaluating $L_3$ at the boolean assignment $(x, y) = (1, 1)$ gives
\begin{equation}
  L_3(1, 1) \;=\; 2 + 2 - 2 - 4 - 2 \;=\; -4 \;<\; 0.
\end{equation}
Furthermore, in Lean~4 (\texttt{pnp\_03\_pairwise\_2sat\_projection\_refuted} and \texttt{pnp\_03\_clausewise\_vs\_global\_existential\_quantifier\_gap}), we formalize the exact quantifier gap on the 8-clause $3\text{-CNF}$ contradiction $\Phi(x_1, x_2, x_3) = \bigwedge_{a,b,c \in \{0,1\}} (x_1^a \lor x_2^b \lor x_3^c)$:
\begin{enumerate}
  \item Under \emph{clause-wise existential projection} $\bigwedge_{i=1}^{8} \exists z_i \in \{0,1\},\, C_i(x_1, x_2, z_i)$ (where each clause chooses its own witness $z_i$), the projected relaxation evaluates to $\mathrm{true}$ for every $(x_1, x_2) \in \{0, 1\}^2$.
  \item Under \emph{global existential projection} $\exists z \in \{0,1\} \bigwedge_{i=1}^{8} C_i(x_1, x_2, z)$ (where a single witness $z$ must simultaneously satisfy all 8 clauses), the projection evaluates to $\mathrm{false}$ for every $(x_1, x_2) \in \{0, 1\}^2$, and $\Phi(x_1, x_2, x_3) = \mathrm{false}$ for all $(x_1, x_2, x_3) \in \{0, 1\}^3$.
\end{enumerate}
\end{counterexample}

\subsection{Refutation of Zero LP Integrality Gap on Odd Cycles (\texorpdfstring{$\mathrm{PNP\text{-}04}$}{PNP-04})}

\begin{candidate}[$\mathrm{PNP\text{-}04}$ --- FALSE]\label{cand:pnp04}
The symmetric Linear Programming relaxation of Vertex Cover / $3\text{-SAT}$ on the odd triangle $K_3$ ($x_1 + x_2 \ge 1, x_2 + x_3 \ge 1, x_3 + x_1 \ge 1$) has zero integrality gap, so the integrality-gap quadratic $G_{\mathrm{LP}}(b) = 16b^2 - 16b + 3 = 4(2b - 1)^2 - 1$ satisfies $G_{\mathrm{LP}}(b) \ge 0$ for all $b \in [0, 1] \cap \Q$.
\end{candidate}

\begin{counterexample}[Refutation of Candidate Lemma~\ref{cand:pnp04}]
At the integral endpoints $b = 0$ and $b = 1$, $G_{\mathrm{LP}}(0) = G_{\mathrm{LP}}(1) = 3 > 0$. However, at the unique symmetric LP fractional optimum $b = \frac{1}{2}$,
\begin{equation}
  G_{\mathrm{LP}}\!\left(\frac{1}{2}\right) \;=\; 16\left(\frac{1}{4}\right) - 16\left(\frac{1}{2}\right) + 3 \;=\; 4 - 8 + 3 \;=\; -1 \;<\; 0.
\end{equation}
In \texttt{lean/PvsNPMathlib.lean} (\texttt{pnp\_04\_k3\_lp\_integrality\_gap\_over\_reals}), we prove the complete LP and IP optimality bounds over $\R$ and $\mathrm{Bool}$:
\begin{enumerate}
  \item Every fractional LP-feasible point $(x_1, x_2, x_3) \in \R^3$ satisfying $x_1+x_2 \ge 1$, $x_2+x_3 \ge 1$, $x_3+x_1 \ge 1$ satisfies $x_1 + x_2 + x_3 \ge \frac{3}{2}$, and $(\frac{1}{2}, \frac{1}{2}, \frac{1}{2})$ achieves $\frac{3}{2}$.
  \item Every integral vertex cover $(x_1, x_2, x_3) \in \{0, 1\}^3$ of $K_3$ satisfies $\sum_{i=1}^{3} x_i \ge 2$, and $(1, 1, 0)$ achieves $2$, establishing the exact integrality gap $\frac{2}{3/2} = \frac{4}{3}$.
\end{enumerate}
\end{counterexample}

\subsection{Refutation of Bilinear Margin Non-Negativity and Karatsuba's Ring Identity (\texorpdfstring{$\mathrm{PNP\text{-}05}$}{PNP-05})}

\begin{candidate}[$\mathrm{PNP\text{-}05}$ --- FALSE]\label{cand:pnp05}
The candidate bilinear margin $K(x, y) = (x + y)^2 - (x - y)^2 - 5xy$ satisfies $K(x, y) \ge 0$ for all $(x, y) \in \R^2$.
\end{candidate}

\begin{proof}
Expanding $K(x, y) = (x^2 + 2xy + y^2) - (x^2 - 2xy + y^2) - 5xy = -xy$ and evaluating at $(x, y) = (1, 1)$ gives
\begin{equation}
  K(1, 1) \;=\; 4 - 0 - 5 \;=\; -1 \;<\; 0.
\end{equation}
Furthermore, in Lean~4 (\texttt{pnp\_05\_karatsuba\_three\_mult\_identity} and \texttt{pnp\_05\_karatsuba\_commRing\_identity}) we prove that in any commutative ring $R$, the bilinear cross-term $a_1 b_0 + a_0 b_1$ satisfies the algebraic cancellation identity
\begin{equation}
  (a_1 + a_0)(b_1 + b_0) - a_1 b_1 - a_0 b_0 \;=\; a_1 b_0 + a_0 b_1,
\end{equation}
expressing the degree-1 convolution cross-term as a linear combination of the three products $z_2 = a_1 b_1$, $z_0 = a_0 b_0$, and $z_3 = (a_1 + a_0)(b_1 + b_0)$. (Note: this formalizes the algebraic ring identity and refutes $K(x,y) \ge 0$; it does not formalize a general Boolean or algebraic circuit lower-bound theorem.)
\end{counterexample}

% ============================================================================
\section{Inductive Tseitin 3-CNF Circuit Transformation (\texorpdfstring{$\mathrm{PNP\text{-}07}$}{PNP-07})}
% ============================================================================

Extending the single-gate gadget of Section~3 ($\mathrm{PNP\text{-}02}$) to arbitrary-depth boolean circuits, we define the inductive syntax $\mathrm{BoolCircuit}$ (`var`, `not`, `and`, `or`) over input variables indexed by $\N$, recursive evaluation $\mathrm{evalCircuit}(\mathrm{env}, C) \in \{0, 1\}$, internal gate count $s = \mathrm{gateCount}(C)$, emitted 3-CNF clause count $\mathrm{tseitinGateClauseCount}(C)$ ($2$ clauses per $\mathrm{not}$ gate and $3$ clauses per $\mathrm{and}$/$\mathrm{or}$ gate), and the recursive wire-consistency predicate $\mathrm{tseitinWireConsistent}(\mathrm{env}, \mathrm{wire}, C)$.

\begin{theorem}[$\mathrm{PNP\text{-}07}$: Structural Induction Equi-Satisfiability and Linear Clause Bound for Arbitrary-Depth Boolean Circuits]\label{thm:pnp07}
For every boolean circuit $C : \mathrm{BoolCircuit}$ with $s = \mathrm{gateCount}(C)$ internal gates and every input environment $\mathrm{env} : \N \to \{0, 1\}$:
\begin{enumerate}
  \item \textbf{(Linear 3-CNF Clause-Count Bound)}: The total number of gate-consistency 3-CNF clauses satisfies $\mathrm{tseitinGateClauseCount}(C) \le 3s$, and including the top-level output unit clause yields at most $3s + 1$ clauses (\texttt{pnp\_07\_tseitin\_clause\_count\_linear\_bound}).
  \item \textbf{(Soundness by Structural Induction)}: Every wire labeling $\mathrm{wire} : \mathrm{BoolCircuit} \to \{0, 1\}$ satisfying $\mathrm{tseitinWireConsistent}(\mathrm{env}, \mathrm{wire}, C)$ satisfies $\mathrm{wire}(C) = \mathrm{evalCircuit}(\mathrm{env}, C)$ (\texttt{pnp\_07\_tseitin\_circuit\_soundness}).
  \item \textbf{(Completeness by Structural Induction)}: The canonical wire assignment $\mathrm{wire}(C') \coloneqq \mathrm{evalCircuit}(\mathrm{env}, C')$ satisfies $\mathrm{tseitinWireConsistent}(\mathrm{env}, \mathrm{evalCircuit}(\mathrm{env}), C)$ (\texttt{pnp\_07\_tseitin\_circuit\_completeness}).
  \item \textbf{(Global Equi-Satisfiability)}: $\mathrm{evalCircuit}(\mathrm{env}, C) = 1$ if and only if there exists a wire assignment $\mathrm{wire} : \mathrm{BoolCircuit} \to \{0, 1\}$ such that $\mathrm{tseitinWireConsistent}(\mathrm{env}, \mathrm{wire}, C)$ and $\mathrm{wire}(C) = 1$ (\texttt{pnp\_07\_tseitin\_circuit\_equisat}).
\end{enumerate}
\end{theorem}

\begin{proof}
We proceed by structural induction on $C$. For base case $C = \mathrm{var}(i)$, $\mathrm{gateCount}(C) = 0 = \mathrm{tseitinGateClauseCount}(C)$ and $\mathrm{tseitinWireConsistent}$ reduces to $\mathrm{wire}(\mathrm{var}(i)) = \mathrm{env}(i) = \mathrm{evalCircuit}(\mathrm{env}, \mathrm{var}(i))$. For inductive steps $C = \mathrm{not}(C_1)$, $\mathrm{and}(C_1, C_2)$, and $\mathrm{or}(C_1, C_2)$, the clause bound follows from $2 \le 3$ and the inductive hypotheses on $C_1, C_2$; soundness follows by substituting $\mathrm{wire}(C_k) = \mathrm{evalCircuit}(\mathrm{env}, C_k)$ into the gate consistency equation; and completeness holds by reflexivity of $\mathrm{evalCircuit}$ at each gate.
\end{proof}

% ============================================================================
\section{Polynomial-Method Degree Lower Bound for 3-Bit Parity (\texorpdfstring{$\mathrm{PNP\text{-}08}$}{PNP-08})}
% ============================================================================

A foundational method in circuit complexity (Razborov 1987; Smolensky 1987) establishes lower bounds by proving that low-degree polynomials cannot represent symmetric functions such as $\mathrm{PARITY}_n$. On $\{0, 1\}^3$, let $\mathrm{paritySign3}(x_1, x_2, x_3) \coloneqq (1 - 2x_1)(1 - 2x_2)(1 - 2x_3) \in \{\pm 1\}$, and define the 3-fold boolean cube finite-difference operator $\Delta_{123} = \Delta_1\Delta_2\Delta_3$:
\begin{equation}
  \Delta_{123}(f) \;\coloneqq\; \sum_{x \in \{0,1\}^3} (-1)^{3 - (x_1 + x_2 + x_3)} f(x_1, x_2, x_3).
\end{equation}

\begin{theorem}[$\mathrm{PNP\text{-}08}$: Unconditional Degree Lower Bound for 3-Bit Parity via Cube Finite Differences]\label{thm:pnp08}
Let $P(x_1, x_2, x_3) = c_0 + \sum_{i=1}^3 c_i x_i + \sum_{i=1}^3 c_{ii} x_i^2 + \sum_{1 \le i < j \le 3} c_{ij} x_i x_j$ be any 3-variable polynomial of total degree at most $2$ over $\R$ (or $\Z$). Then:
\begin{enumerate}
  \item \textbf{(Annihilation of Degree-$\le 2$ Polynomials)}: $\Delta_{123}(P) = 0$ identically for all coefficients $c_0, \dots, c_{13}$ (\texttt{pnp\_08\_cubeFiniteDiff3\_degAtMost2\_eq\_zero}).
  \item \textbf{(Non-Vanishing Cube Derivative of Parity)}: $\Delta_{123}(\mathrm{paritySign3}) = -8 \ne 0$ (\texttt{pnp\_08\_cubeFiniteDiff3\_paritySign3\_eq\_neg\_eight}).
  \item \textbf{(Strict Degree Lower Bound)}: No polynomial $P$ of total degree at most $2$ over $\R$ (or $\Z$) can agree with $\mathrm{paritySign3}$ on all $8$ vertices of $\{0, 1\}^3$; hence any polynomial computing 3-bit Parity on $\{0, 1\}^3$ has degree at least $3$ (\texttt{pnp\_08\_no\_deg2\_polynomial\_computes\_parity3}).
\end{enumerate}
\end{theorem}

\begin{proof}
Each difference operator $\Delta_i f(x) = f(x|_{x_i=1}) - f(x|_{x_i=0})$ annihilates any monomial not containing $x_i$. Since every monomial in $P$ has total degree $\le 2$, it omits at least one variable $x_i \in \{x_1, x_2, x_3\}$, so $\Delta_{123}(P) = 0$ by ring expansion. Conversely, $\mathrm{paritySign3}(x_1,x_2,x_3) = (-1)^{x_1+x_2+x_3}$ has the same sign as $(-1)^{3-(x_1+x_2+x_3)} = -(-1)^{x_1+x_2+x_3}$ on all $8$ vertices of $\{0,1\}^3$, giving $\Delta_{123}(\mathrm{paritySign3}) = 8 \times (-1) = -8 \ne 0$. If $P(x) = \mathrm{paritySign3}(x)$ for all $x \in \{0,1\}^3$, then $0 = \Delta_{123}(P) = \Delta_{123}(\mathrm{paritySign3}) = -8$, a contradiction.
\end{proof}

% ============================================================================
\section{Summary of Dual-Kernel Formal Verification}
% ============================================================================

\begin{table}[htbp]
\centering
\small
\begin{tabular}{@{}llll@{}}
\toprule
\textbf{Unit} & \textbf{Subject} & \textbf{Formal Verdict} & \textbf{Lean~4 Theorem / Certificate} \\
\midrule
$\mathrm{PNP\text{-}01}$ & Single-Variable Quartic Penalty & \textbf{PROVEN\_TRUE} & \texttt{pnp\_01\_cookLevinPenalty\_sos\_and\_zero\_locus} \\
                         & $E_1(x) = (x^2-x)^2 + (x-1)^2 \ge 0$ & ($\mathrm{SOS}$ Identity) & $E_1(x) = 0 \iff x = 1$ \\
$\mathrm{PNP\text{-}02}$ & Tseitin $\mathrm{AND}$-Gate \& Slack & \textbf{PROVEN\_TRUE} & \texttt{pnp\_02\_tseitin\_and\_gadget\_equisat} \\
                         & $\mathcal{T}_{\land}(x,y,z) \iff z = x\land y$ & ($\mathrm{SOS}$ + `Bool`) & + \texttt{pnp\_02\_tseitinRealSlack\_sos} \\
$\mathrm{PNP\text{-}03}$ & Pairwise $2\text{-SAT}$ \& Quantifier Gap & \textbf{REFUTED\_FALSE} & \texttt{pnp\_03\_pairwise\_2sat\_projection\_refuted} \\
                         & $L_3(1,1) = -4 < 0$ & ($\forall i \exists z_i$ vs $\exists z \forall i$) & + \texttt{pnp\_03\_clausewise\_vs\_global\_existential\_quantifier\_gap} \\
$\mathrm{PNP\text{-}04}$ & Odd-Cycle $K_3$ LP/IP Optimality & \textbf{REFUTED\_FALSE} & \texttt{pnp\_04\_k3\_lp\_integrality\_gap\_over\_reals} \\
                         & $\mathrm{LP} = \frac{3}{2}$ vs.\ $\mathrm{IP} = 2$ & (Exact Gap $\frac{4}{3}$) & + \texttt{pnp\_04\_k3\_lp\_integrality\_gap\_refuted} \\
$\mathrm{PNP\text{-}05}$ & Bilinear Margin \& Karatsuba Identity & \textbf{REFUTED\_FALSE} & \texttt{pnp\_05\_non\_monotone\_cancellation\_refuted} \\
                         & $K(x,y) = -xy,\; K(1,1)=-1<0$ & (`[CommRing R]`) & + \texttt{pnp\_05\_karatsuba\_commRing\_identity} \\
$\mathrm{PNP\text{-}06}$ & Spectral Discrepancy \& Diagonalization & \textbf{PROVEN\_TRUE} & \texttt{pnp\_06\_baker\_gill\_solovay\_diagonal\_differs} \\
                         & $D_{\mathrm{spec}}(x,y) = (x-y)^2 \ge 0$ & (Pointwise Diag.) & + \texttt{pnp\_06\_no\_machine\_decides\_diagonalLanguage} \\
$\mathrm{PNP\text{-}07}$ & Inductive Tseitin 3-CNF Circuit & \textbf{PROVEN\_TRUE} & \texttt{pnp\_07\_tseitin\_circuit\_equisat} \\
                         & $\le 3s+1$ Clauses \& Equi-Satisfiable & (`BoolCircuit`) & + \texttt{pnp\_07\_tseitin\_clause\_count\_linear\_bound} \\
$\mathrm{PNP\text{-}08}$ & Polynomial Degree Bound ($\mathrm{Parity}_3$) & \textbf{PROVEN\_TRUE} & \texttt{pnp\_08\_no\_deg2\_polynomial\_computes\_parity3} \\
                         & $\Delta_{123}(P_{\le 2})=0 \ne -8=\Delta_{123}(\mathrm{Parity}_3)$ & (Cube Finite Diff.) & + \texttt{pnp\_08\_cubeFiniteDiff3\_degAtMost2\_eq\_zero} \\
\bottomrule
\end{tabular}
\caption{Summary of Dual Symbolic and Lean~4 Kernel Adjudications.}
\label{tab:pnp-g3-summary}
\end{table}

\end{document}
